% Level 1, Workbook chapter 10 -- Random Art and My Own Commands   (input by student-workbook.tex)
\chapterweek{10}{Random Art and My Own Commands}{\texttt{random} makes surprises; \texttt{def} makes your own command, just like My Blocks.}

\section{What you will learn}
\begin{itemize}[leftmargin=1.2em,itemsep=1pt]
  \item How to make your own command with \texttt{def}, with input slots called \textbf{parameters}.
  \item The difference between \textbf{defining} a command and \textbf{calling} (using) it.
  \item How \texttt{random} makes a new picture every time you press Run.
  \item How to plan your final project: \textbf{My Interactive Art Show}.
\end{itemize}

\section{The big idea: make your own command}
In Scratch you made My Blocks. Python has them too! You write \texttt{def}, a name, and the input slots in brackets.
\begin{center}
\begin{tabular}{@{}m{6.2cm} m{0.6cm} m{7.2cm}@{}}
\begin{scratch}[scale=0.8]
\initmoreblocks{define \namemoreblocks{flower} \ovalmoreblocks{x} \ovalmoreblocks{y} \ovalmoreblocks{colour}}
\blockpen{pen up}
\blockmove{go to x: \ovalmoreblocks{x} y: \ovalmoreblocks{y}}
\blockpen{pen down}
\end{scratch}
& → &
\begin{pycode}
def flower(x, y, colour):
    t.penup()
    t.goto(x, y)
    t.pendown()
    t.color(colour)
    ...
\end{pycode}
\end{tabular}
\end{center}
\textbf{Defining} writes the recipe. Nothing is drawn yet! \textbf{Calling} uses the recipe. The values in the brackets go into the slots, in order:
\begin{center}
\begin{tikzpicture}[font=\small]
  \node[font=\ttfamily\large] (call) at (0,0) {flower(-200, 0, "\#E69F00")};
  \node[draw,thick,rounded corners=4pt,fill=kcyellow!30,minimum width=3.2cm,minimum height=1.5cm,align=center] (m) at (6.6,0) {\textbf{flower machine}\\\texttt{x = -200}\\\texttt{y = 0}\\\texttt{colour = orange}};
  \draw[-{Stealth},very thick] (call) -- (m);
  \begin{scope}[shift={(10.4,0)}]
    \foreach \a in {0,60,...,300}
      \fill[kcorange,draw=black] (\a:0.42) circle (0.32);
    \fill[kcyellow,draw=black] (0,0) circle (0.22);
  \end{scope}
  \draw[-{Stealth},very thick] (m) -- (9.5,0);
  \node[font=\small] at (10.4,-1.05) {a flower at (-200, 0)};
\end{tikzpicture}
\end{center}
\begin{tip}[Remember]
Define first (at the top), call later (at the bottom). Calls start at the left edge --- not pushed in --- or they become part of the recipe.
\end{tip}

\begin{exercise}[Flower Stamp]
Folder: \texttt{week10/ex-flower-stamp}. Run the starter: nothing is drawn! The recipe is there, but nobody calls it.
\begin{enumerate}[leftmargin=1.4em,itemsep=1pt]
  \item At the bottom, call it three times: \texttt{flower(-200, 0)}, \texttt{flower(0, 0)}, \texttt{flower(200, 0)}.
  \item Add a third parameter: \texttt{def flower(x, y, colour):} and use \texttt{t.color(colour)}.
  \item Give each call a colour: \texttt{flower(-200, 0, PALETTE[0])}, and so on.
\end{enumerate}
Done looks like: three flowers in different colours, each with a yellow middle. Then run \texttt{check\_my\_work.py}.
\end{exercise}

\section{Surprise! The random module}
\begin{center}
\renewcommand{\arraystretch}{1.5}
\begin{tabularx}{\linewidth}{@{}p{6.9cm}>{\raggedright\arraybackslash}X@{}}
\toprule
\texttt{random.randint(1, 6)} & Like rolling a dice: a whole number from 1 to 6.\\
\texttt{random.randint(-350, 350)} & A random spot from left (-350) to right (350).\\
\texttt{random.choice(PALETTE)} & Like picking a crayon from a bag with your eyes closed.\\
\bottomrule
\end{tabularx}
\end{center}
Put \texttt{import random} at the top of your program, next to \texttt{import turtle}.

\begin{exercise}[Random Bubble Art]
Folder: \texttt{week10/ex-random-art}. The starter draws 30 bubbles\ldots{} all in the same place! Inside the loop, change the fixed numbers to random ones: \texttt{x}, \texttt{y}, \texttt{size} with \texttt{random.randint}, and \texttt{colour} with \texttt{random.choice(PALETTE)}. Press Run three times. Is it different each time?
\end{exercise}

\section{Your final project: My Interactive Art Show}
For weeks 10 to 12 you will make your own program for an art show. It says hello, asks the viewer one or two questions, and draws a picture that changes because of the answers. It must use these \textbf{six ingredients}:
\begin{checklist}
  \item a question with \texttt{input}
  \item a variable that remembers an answer or a number
  \item an \texttt{if} that chooses what to draw
  \item a \texttt{for} loop that repeats something
  \item \texttt{random} for a surprise
  \item a function you made with \texttt{def}
\end{checklist}
\textbf{Ideas:} a fish party under the sea, a garden of flowers, a city skyline at night, space with planets, a birthday card, a pattern machine. Write three ideas you like:
\answerlines{3}
Then fill in the \textbf{plan sheet} in your Project Brief (Milestone M1): a sketch, your questions, and your function's name and inputs.

\begin{tip}[Be an original artist]
Borrowing an idea is fine when you make it your own and say thank you. Write a comment like \texttt{\# Flower idea remixed from Exercise 10.1}. Questions in your project ask about choices --- colours, numbers, day or night --- never about who someone is.
\end{tip}

\begin{vocab}
  \vocabword{Function}{A command you make yourself, with a name. Like a My Block.}
  \vocabword{def}{The Python word that starts a function's recipe: \texttt{def flower(x, y, colour):}.}
  \vocabword{Parameter}{An input slot of a function, like \texttt{x}, \texttt{y} or \texttt{colour}.}
  \vocabword{Call}{Using a function: \texttt{flower(0, 0, "\#CC79A7")}.}
  \vocabword{random}{A Python toolbox for surprises: \texttt{randint} and \texttt{choice}.}
  \vocabword{Remix}{Making something new from someone's idea, and giving them credit.}
\end{vocab}

\begin{reflection}
What command would you like to invent for your project? What would its input slots be?
\answerlines{2}
Why do you think a function is useful when you need ten flowers?
\answerlines{2}
\end{reflection}
